\documentclass[12pt,a4paper]{article}

% Paquetes básicos
\usepackage[utf8]{inputenc}
\usepackage[spanish]{babel}
\usepackage{geometry}
\geometry{a4paper, margin=2.5cm}
\usepackage{hyperref}
\usepackage{titlesec}
\usepackage{enumitem}

% Configuración de enlaces
\hypersetup{
    colorlinks=true,
    linkcolor=blue,
    filecolor=magenta,      
    urlcolor=cyan,
}

% Título y autores
\title{\textbf{Documentación: La vigilancia tecnológica en la gestión de proyectos de I+D+i}}
\author{Basado en el texto de: Javier Muñoz Durán, María Marín Martínez y José Vallejo Triano}
\date{\today}

\begin{document}

\maketitle

\begin{abstract}
Este documento resume los procesos y herramientas fundamentales para la implantación de un sistema de Vigilancia Tecnológica (VT) en organizaciones orientadas a la I+D+i, tomando como referencia la norma UNE 166006:2006 EX.
\end{abstract}

\section{Introducción a la Vigilancia Tecnológica (VT)}
La vigilancia tecnológica es una forma organizada, selectiva y permanente de captar información del exterior, analizarla y convertirla en conocimiento[cite: 79]. Esto permite tomar decisiones con menor riesgo y anticiparse a los cambios en el entorno de la innovación[cite: 79].

\section{Fases del Proceso de Vigilancia Tecnológica}
Para integrar la VT en el sistema de gestión, el artículo propone dividirla en los siguientes procesos clave[cite: 84, 85]:

\subsection{1. Identificación de necesidades}
Las necesidades se determinan mediante un autodiagnóstico de la organización[cite: 88]. A partir de este, se obtienen los \textbf{factores críticos de vigilancia (FCVs)}, que son cuestiones externas cuya evolución afecta la competitividad[cite: 89, 90].

\subsection{2. Identificación de fuentes}
Una vez detectados los FCVs, es necesario clasificar las fuentes de información[cite: 99, 107]:
\begin{itemize}
    \item \textbf{Fuentes informales:} Conversaciones con clientes, proveedores e investigadores. Son de alto valor pero difíciles de estructurar[cite: 108, 122].
    \item \textbf{Fuentes formales:} Bases de datos o servicios de internet. Su información está homogeneizada, lo que permite automatizar su explotación[cite: 123, 124].
\end{itemize}

\subsection{3. Medios de acceso}
Se refieren a las herramientas utilizadas para obtener la información de las fuentes[cite: 131]. Es importante destacar la existencia de la ``web profunda'', bases de datos accesibles desde internet que los buscadores convencionales no pueden indexar[cite: 216, 217].

\subsection{4. Búsqueda}
La búsqueda es un proceso iterativo que requiere chequear constantemente si los resultados se corresponden con lo esperado[cite: 252]. Exige definir una estrategia previa con términos controlados y fuentes pertinentes[cite: 253, 254].

\subsection{5. Puesta en valor de la información}
Consiste en validar y analizar la información para la toma de decisiones[cite: 244]. Se pueden utilizar técnicas avanzadas como la minería de datos aplicada a fuentes estructuradas, realizando recuentos y análisis de coocurrencias[cite: 262, 263].

\subsection{6. Difusión de la información}
Es el último paso del proceso, donde se realiza la difusión selectiva de la información generada hacia el usuario según su perfil, generalmente utilizando el correo electrónico.

\section{Las Patentes como Fuente Estratégica}
El artículo resalta el análisis de las patentes debido a la exclusividad y estructuración de sus contenidos[cite: 63, 229]. Se estima que más del 80\% de la información técnica actual está contenida en estos documentos y no se publica en ningún otro medio[cite: 63, 226]. Analizarlas permite observar las líneas de desarrollo de la competencia e identificar nuevas oportunidades[cite: 228].

\section{Conclusión}
La correcta implantación de un sistema de VT, apoyado en herramientas de software y metodologías estructuradas, permite a las organizaciones innovadoras mejorar sus inversiones, elegir líneas de I+D óptimas y anticiparse a su competencia[cite: 287, 288, 289].

\end{document}